\documentclass[11pt,letterpaper]{article}
\usepackage[technical-reference]{cornell-notes}
\usepackage{needspace}

\title{Clause 4: General Principles - Cornell Notes}
\author{}
\date{}
\setDocTitle{Clause 4: General Principles}
\setDocSubtitle{C++ 2024 Cornell Notes}
\setDocOwner{C++ 2024 Cornell Notes}
\setDocAuthor{}
\setDocDate{}
\setCornellCollection{C++ 2024 Cornell Notes}
\setCornellUnitType{Clause}
\setCornellUnitNumber{4}
\setCornellUnitTitle{General Principles}
\setCornellCompanionLabel{C++ Technical Reference}
\hypersetup{pdftitle={Clause 4: General Principles - Cornell Notes},pdfsubject={C++ technical study notes},pdfauthor={}}

\begin{document}
\maketitle
\begin{CornellReferenceOrientationBox}
\textbf{Purpose.} Explains implementation compliance, the abstract-machine model, document organization, and syntax notation. It distinguishes diagnosable violations, no-diagnostic-required cases, undefined behavior, hosted and freestanding obligations, and the meaning of grammar notation.
\par\textbf{Role.} Supplies the interpretive rules used to read every later core-language and library requirement.
\par\textbf{Principal dependencies.} Uses the behavior vocabulary from Clause 3 and governs the interpretation of Clauses 5 through 33.
\par\textbf{Primary audience.} programmers assessing portability and implementers assessing conformance.
\par\textbf{Major subject groups.} Implementation compliance; Structure of this document; Syntax notation.
\end{CornellReferenceOrientationBox}
\section{Learning Objectives}
\begin{itemize}[label=\textcolor{CornellReferenceTeal}{$\blacktriangleright$}]
\item Define the governing ideas of Implementation compliance and state the consequences for a program or implementation.
\item Identify the governing ideas of Structure of this document and state the consequences for a program or implementation.
\item Distinguish the governing ideas of Syntax notation and state the consequences for a program or implementation.
\item Explain the governing ideas of General and state the consequences for a program or implementation.
\item Trace the governing ideas of Abstract machine and state the consequences for a program or implementation.
\end{itemize}
\section{Normative Structure Map}
\small\begin{longtable}{@{}>{\RaggedRight\arraybackslash}p{0.13\textwidth}>{\RaggedRight\arraybackslash}p{0.27\textwidth}>{\RaggedRight\arraybackslash}p{0.19\textwidth}>{\RaggedRight\arraybackslash}p{0.31\textwidth}@{}}
\toprule\textbf{Subclause} & \textbf{Subject} & \textbf{Rule category} & \textbf{Key dependencies / entities}\\\midrule\endfirsthead
\toprule\textbf{Subclause} & \textbf{Subject} & \textbf{Rule category} & \textbf{Key dependencies / entities}\\\midrule\endhead
\bottomrule\endfoot
4.1 & Implementation compliance & core-language semantics & General, Abstract machine\\
4.2 & Structure of this document & core-language semantics & Uses the behavior vocabulary from Clause 3 and governs the interpretation of Clauses 5 through 33.\\
4.3 & Syntax notation & syntax / lexical rules & Uses the behavior vocabulary from Clause 3 and governs the interpretation of Clauses 5 through 33.\\
\end{longtable}\normalsize
\begin{CornellReferenceNormativeBox}A heading maps the clause; it does not replace the detailed conditions. For each operation, read requirements, effects, result, exceptions, complexity, remarks, and notes with their stated normative force.\end{CornellReferenceNormativeBox}
\subsection*{Rule Classification Guide}
\small\begin{longtable}{@{}>{\RaggedRight\arraybackslash}p{0.25\textwidth}>{\RaggedRight\arraybackslash}p{0.67\textwidth}@{}}\toprule\textbf{Label or category} & \textbf{How to read it}\\\midrule\endfirsthead\toprule\textbf{Label or category} & \textbf{How to read it (continued)}\\\midrule\endhead\bottomrule\endfoot
Well-formed / ill-formed & Formation is governed by syntax and semantic rules; an ill-formed program does not become well-formed because one implementation accepts it.\\
Diagnosable rule & A violation requires at least one diagnostic unless the wording places the case outside the diagnosable set.\\
No diagnostic required & The program can be ill-formed while the implementation has no diagnostic obligation.\\
Undefined behavior & No execution requirements are imposed; translation success does not establish safety or portability.\\
Unspecified behavior & One of the permitted outcomes may be chosen without documentation.\\
Implementation-defined behavior & The implementation chooses and documents the behavior.\\
Conditionally supported & The implementation may omit support but documents unsupported constructs.\\
\end{longtable}\normalsize
\section{Cornell Cue-and-Notes}
\Needspace{8\baselineskip}
\subsection*{4.1 Implementation compliance}
\begin{longtable}{@{}>{\RaggedRight\arraybackslash\columncolor{CornellReferenceCueGray}}p{0.28\textwidth}>{\RaggedRight\arraybackslash}p{0.64\textwidth}@{}}
\rowcolor{CornellReferenceNavy}\color{white}\textbf{Cue / recall prompt} & \color{white}\textbf{Notes}\\\endfirsthead
\rowcolor{CornellReferenceNavy}\color{white}\textbf{Cue / recall prompt} & \color{white}\textbf{Notes (continued)}\\\endhead
\arrayrulecolor{CornellReferenceLineGray}
\textbf{What rule applies to General?}\\[-1mm]{\scriptsize\CornellClauseRef{4.1.1}} & Frames the terminology, applicability, and relationships needed to interpret General; detailed rules remain in the following subclauses.\\\hline
\textbf{What rule applies to Abstract machine?}\\[-1mm]{\scriptsize\CornellClauseRef{4.1.2}} & Defines the core-language rules for Abstract machine, including formation, semantic effect, interactions with type and lookup, and any ill-formed or implementation-dependent cases.\\\hline
\end{longtable}
\begin{CornellReferenceSummaryBox}Implementation compliance supplies the core-language semantics portion of this clause. The key review move is to connect its local wording to the clause-wide model, then classify each condition by normative force before relying on it.\end{CornellReferenceSummaryBox}
\Needspace{8\baselineskip}
\subsection*{4.2 Structure of this document}
\begin{longtable}{@{}>{\RaggedRight\arraybackslash\columncolor{CornellReferenceCueGray}}p{0.28\textwidth}>{\RaggedRight\arraybackslash}p{0.64\textwidth}@{}}
\rowcolor{CornellReferenceNavy}\color{white}\textbf{Cue / recall prompt} & \color{white}\textbf{Notes}\\\endfirsthead
\rowcolor{CornellReferenceNavy}\color{white}\textbf{Cue / recall prompt} & \color{white}\textbf{Notes (continued)}\\\endhead
\arrayrulecolor{CornellReferenceLineGray}
\textbf{What rule applies to Structure of this document?}\\[-1mm]{\scriptsize\CornellClauseRef{4.2}} & Defines the core-language rules for Structure of this document, including formation, semantic effect, interactions with type and lookup, and any ill-formed or implementation-dependent cases.\\\hline
\end{longtable}
\begin{CornellReferenceSummaryBox}Structure of this document supplies the core-language semantics portion of this clause. The key review move is to connect its local wording to the clause-wide model, then classify each condition by normative force before relying on it.\end{CornellReferenceSummaryBox}
\Needspace{8\baselineskip}
\subsection*{4.3 Syntax notation}
\begin{longtable}{@{}>{\RaggedRight\arraybackslash\columncolor{CornellReferenceCueGray}}p{0.28\textwidth}>{\RaggedRight\arraybackslash}p{0.64\textwidth}@{}}
\rowcolor{CornellReferenceNavy}\color{white}\textbf{Cue / recall prompt} & \color{white}\textbf{Notes}\\\endfirsthead
\rowcolor{CornellReferenceNavy}\color{white}\textbf{Cue / recall prompt} & \color{white}\textbf{Notes (continued)}\\\endhead
\arrayrulecolor{CornellReferenceLineGray}
\textbf{What rule applies to Syntax notation?}\\[-1mm]{\scriptsize\CornellClauseRef{4.3}} & Defines the recognized forms for Syntax notation; syntactic acceptance does not replace the semantic constraints stated with the production.\\\hline
\end{longtable}
\begin{CornellReferenceSummaryBox}Syntax notation supplies the syntax / lexical rules portion of this clause. The key review move is to connect its local wording to the clause-wide model, then classify each condition by normative force before relying on it.\end{CornellReferenceSummaryBox}
\section{Language-Lawyer and Library-Usage Pitfalls}
\begin{CornellReferencePitfallBox}\begin{itemize}[label=\textcolor{CornellReferenceAmber}{$\blacktriangleright$}]
\item Treating successful translation as proof that a program has defined behavior.
\item Expecting a diagnostic for undefined behavior or an explicit no-diagnostic-required case.
\item Assuming a freestanding implementation has every hosted library facility.
\item Reading a note or example as a mandatory rule.
\item Treating a note, example, or recommended practice as though it imposed a mandatory program rule.
\end{itemize}\end{CornellReferencePitfallBox}
\section{Programmer and Implementation Checklist}
\begin{CornellReferenceChecklistBox}\begin{itemize}[label=$\square$]
\item Classify a violation before deciding whether a diagnostic is required.
\item Identify the observable behavior constrained by the abstract machine.
\item Check hosted or freestanding applicability.
\item Separate normative requirements from notes, examples, and recommended practice.
\item Interpret grammar productions together with their semantic constraints.
\item For \CornellClauseRef{4.1}, verify implementation compliance against the precise rule category and all cited dependencies.
\item For \CornellClauseRef{4.2}, verify structure of this document against the precise rule category and all cited dependencies.
\item For \CornellClauseRef{4.3}, verify syntax notation against the precise rule category and all cited dependencies.
\end{itemize}\end{CornellReferenceChecklistBox}
\section{Clause Synthesis}
\begin{CornellReferenceOrientationBox}Explains implementation compliance, the abstract-machine model, document organization, and syntax notation. It distinguishes diagnosable violations, no-diagnostic-required cases, undefined behavior, hosted and freestanding obligations, and the meaning of grammar notation. Supplies the interpretive rules used to read every later core-language and library requirement. A sound review distinguishes syntax and participation from runtime preconditions, keeps notes and examples non-mandatory, and follows cross-clause dependencies before making a portability claim.\end{CornellReferenceOrientationBox}
\section{Self-Test Questions}
\begin{enumerate}
\item What role does Implementation compliance play in the clause's overall model?
\item Which normative distinctions must be kept separate when applying Structure of this document?
\item How would you audit a program or implementation that depends on Syntax notation?
\item Scenario: a program relies on an unstated implementation behavior while using General. What must be checked before calling the program portable?
\item Scenario: a program relies on an unstated implementation behavior while using Abstract machine. What must be checked before calling the program portable?
\end{enumerate}
\clearpage\section{Answer Key}
\begin{enumerate}
\item Implementation compliance is one part of the clause's core-language semantics model. Its role is understood by connecting the local rule in 4.1 to the clause orientation and its dependent concepts.
\item Keep formation and well-formedness separate from runtime preconditions and effects. Also distinguish mandatory wording from notes, implementation choice from undefined behavior, and interface declarations from detailed contracts; see 4.2.
\item Audit Syntax notation by locating the controlling wording in 4.3, listing inputs and type constraints, proving any preconditions, recording effects and failure behavior, and checking lifetime, invalidation, complexity, and synchronization where applicable.
\item First identify the exact requirement in 4.1.1, then classify the relied-on behavior as required, unspecified, implementation-defined, conditionally supported, or undefined. Portability requires satisfying the program obligations and avoiding dependence on a choice not guaranteed in the target environments.
\item First identify the exact requirement in 4.1.2, then classify the relied-on behavior as required, unspecified, implementation-defined, conditionally supported, or undefined. Portability requires satisfying the program obligations and avoiding dependence on a choice not guaranteed in the target environments.
\end{enumerate}
\section{Key-Term Recap}
\begin{longtable}{@{}>{\RaggedRight\arraybackslash}p{0.28\textwidth}>{\RaggedRight\arraybackslash}p{0.64\textwidth}@{}}\toprule\textbf{Term} & \textbf{Working meaning}\\\midrule\endfirsthead\toprule\textbf{Term} & \textbf{Working meaning (continued)}\\\midrule\endhead\bottomrule\endfoot
diagnosable rule & A syntax or semantic rule whose violation requires at least one diagnostic unless explicitly excluded.\\
abstract machine & The semantic model against which observable program behavior is specified.\\
hosted implementation & An implementation providing the complete hosted facility set.\\
freestanding implementation & An implementation intended to operate without the usual operating-system environment and providing the required subset.\\
grammar notation & The formal conventions used to express syntactic alternatives and productions.\\
\end{longtable}
\CornellTakeaway{Conformance analysis begins by classifying the rule, the implementation model, and the normative force of the wording.}
\end{document}
